% GENERATED FILE. Edit canonical lessons, then run npm run generate:books.
% canonical-source-hash: fnv1a64:bbb7c928fcc4a061
% canonical-lessons: RU-C36-children, RU-C36-child, RU-C36-aunt, RU-C36-grandson, RU-C36-granddaughter

\chapter{Children, a Child, Aunt, Grandchildren}
\label{ch:ru-children-a-child-aunt-grandchildren}
\hlchaptermodality{36}

\begin{chapteropening}
\textbf{By the end of this chapter:} \emph{I can name children, a child, an aunt, a grandson and a granddaughter.}

Complete the last lesson of chapter 36: i can name children, a child, an aunt, a grandson and a granddaughter.
\end{chapteropening}

\section[déti]{\ru{дети} (déti) — children}

\label{lesson:RU-C36-children}



\begin{quote}
Before the new one: say the Russian for a husband, then the Russian for a wife.
\end{quote}

\subsection*{You'll want to know: \ru{дети}}
\textbf{\ru{дети}} (\emph{déti}) — \textquotedblleft{}children\textquotedblright{}.

\begin{grammarlens}[title={What you've built}]
One more word for children, a child, aunt, grandchildren.
\end{grammarlens}

\subsection*{Your turn}
\begin{itemize}
\raggedright
  \item \emph{Say it:} \emph{déti}
  \item \emph{Say it:} \emph{déti}, once more
  \item \emph{Say it:} say \emph{déti}
  \item \emph{Recall:} say the Russian for a husband, then the Russian for a wife, then say \emph{déti} again
\end{itemize}

\begin{tcolorbox}[breakable,colback=teal!4,colframe=teal!35!black,title={Before you move on}]
What does \ru{дети} mean? (\textquotedblleft{}Children\textquotedblright{}.) Say it once more.
\end{tcolorbox}
\section[rebyónok]{\ru{ребёнок} (rebyónok) — a child}

\label{lesson:RU-C36-child}



\begin{quote}
Before the new one: say the Russian for a wife, then the Russian for children.
\end{quote}

\subsection*{You'll want to know: \ru{ребёнок}}
\textbf{\ru{ребёнок}} (\emph{rebyónok}) — \textquotedblleft{}a child\textquotedblright{}. One child; \textbf{\ru{дети}} is the plural.

\begin{grammarlens}[title={What you've built}]
One more word for children, a child, aunt, grandchildren.
\end{grammarlens}

\subsection*{Your turn}
\begin{itemize}
\raggedright
  \item \emph{Say it:} \emph{rebyónok}
  \item \emph{Say it:} \emph{rebyónok}, once more
  \item \emph{Say it:} say \emph{rebyónok}
  \item \emph{Recall:} say the Russian for a wife, then the Russian for children, then say \emph{rebyónok} again
\end{itemize}

\begin{tcolorbox}[breakable,colback=teal!4,colframe=teal!35!black,title={Before you move on}]
What does \ru{ребёнок} mean? (\textquotedblleft{}A child\textquotedblright{}.) Say it once more.
\end{tcolorbox}
\section[tyótya]{\ru{тётя} (tyótya) — an aunt}

\label{lesson:RU-C36-aunt}



\begin{quote}
Before the new one: say the Russian for children, then the Russian for a child.
\end{quote}

\subsection*{You'll want to know: \ru{тётя}}
\textbf{\ru{тётя}} (\emph{tyótya}) — \textquotedblleft{}an aunt\textquotedblright{}.

\begin{grammarlens}[title={What you've built}]
One more word for children, a child, aunt, grandchildren.
\end{grammarlens}

\subsection*{Your turn}
\begin{itemize}
\raggedright
  \item \emph{Say it:} \emph{tyótya}
  \item \emph{Say it:} \emph{tyótya}, once more
  \item \emph{Say it:} say \emph{tyótya}
  \item \emph{Recall:} say the Russian for children, then the Russian for a child, then say \emph{tyótya} again
\end{itemize}

\begin{tcolorbox}[breakable,colback=teal!4,colframe=teal!35!black,title={Before you move on}]
What does \ru{тётя} mean? (\textquotedblleft{}An aunt\textquotedblright{}.) Say it once more.
\end{tcolorbox}
\section[vnuk]{\ru{внук} (vnuk) — a grandson}

\label{lesson:RU-C36-grandson}



\begin{quote}
Before the new one: say the Russian for a child, then the Russian for an aunt.
\end{quote}

\subsection*{You'll want to know: \ru{внук}}
\textbf{\ru{внук}} (\emph{vnuk}) — \textquotedblleft{}a grandson\textquotedblright{}.

\begin{grammarlens}[title={What you've built}]
One more word for children, a child, aunt, grandchildren.
\end{grammarlens}

\subsection*{Your turn}
\begin{itemize}
\raggedright
  \item \emph{Say it:} \emph{vnuk}
  \item \emph{Say it:} \emph{vnuk}, once more
  \item \emph{Say it:} say \emph{vnuk}
  \item \emph{Recall:} say the Russian for a child, then the Russian for an aunt, then say \emph{vnuk} again
\end{itemize}

\begin{tcolorbox}[breakable,colback=teal!4,colframe=teal!35!black,title={Before you move on}]
What does \ru{внук} mean? (\textquotedblleft{}A grandson\textquotedblright{}.) Say it once more.
\end{tcolorbox}
\section[vnúchka]{\ru{внучка} (vnúchka) — a granddaughter}

\label{lesson:RU-C36-granddaughter}



\begin{quote}
Before the new one: say the Russian for an aunt, then the Russian for a grandson.
\end{quote}

\subsection*{You'll want to know: \ru{внучка}}
\textbf{\ru{внучка}} (\emph{vnúchka}) — \textquotedblleft{}a granddaughter\textquotedblright{}.

\begin{grammarlens}[title={What you've built}]
One more word for children, a child, aunt, grandchildren.
\end{grammarlens}

\subsection*{Your turn}
\begin{itemize}
\raggedright
  \item \emph{Say it:} \emph{vnúchka}
  \item \emph{Say it:} \emph{vnúchka}, once more
  \item \emph{Say it:} say \emph{vnúchka}
  \item \emph{Recall:} say the Russian for an aunt, then the Russian for a grandson, then say \emph{vnúchka} again
\end{itemize}

\begin{tcolorbox}[breakable,colback=teal!4,colframe=teal!35!black,title={Before you move on}]
What does \ru{внучка} mean? (\textquotedblleft{}A granddaughter\textquotedblright{}.) Say it once more.
\end{tcolorbox}
